\section{The marked observable at the terminal level}
\label{sec:terminal-marked-construction}
The terminal vacuum construction must also carry each local field and
its forced profile. We give this construction before propagating the
field down the selected history. All activities, supports and predicates
in this section belong to that history.

\subsection{Reserved addresses and the missing part of the field}
Fix a grade $d$, a cutoff $N$ and a period $p$. Write $R$ for the
terminal representation, $F$ for its grade-$d$ coefficient frame,
$S$ for its sound local section, $o=0$ for its origin, $Y$ for its
footprint and $0<\theta\le1$ for its threshold. The forced activities are
$K_i$, with supports $X_i$ and nonempty inventories $I_i\subseteq X_i$.
Their eligibility requires both the prescribed flags on $X_i$ and the
original additional predicate. A compatible collection has pairwise
disjoint supports, and its inventories cover exactly the flagged sites.

The activity dictionary used in the startup and return constructions
has, for every pair $I\subseteq X$ with $I\ne\varnothing$, an infinite
sequence of reserved addresses. The address indexed by $(X,I,j)$ has
support $X$, inventory $I$, and zero activity. Its additional predicate
is precisely the condition that the flagged sites in $X$ equal $I$.
Here is why these properties survive the construction. The dictionary
embeds the old hard labels in one set of coordinates and reserves a
disjoint set for every shape $(X,I)$. The new connected objects encode
finite sets of these coordinates. A reserved singleton cannot be an
image of an old object: that would put its reserved coordinate in the
old coordinate embedding. Its coefficient therefore remains zero.
Reading a singleton recovers exactly its support and inventory; its
eligibility test is the original test for that one shape. This is the
same operation in startup and every return. In particular it supplies
the reserved addresses in the actual selected terminal representation.

We use these addresses to encode a finite compatible nonempty family
$b$ of old forced labels whose supports all meet $Y$. Put
\[
 T_b=Y\cup\bigcup_{i\in b}X_i,\qquad
 I_b=\bigcup_{i\in b}I_i.
\]
Choose the reserved address $s(b)$ for shape $(T_b,I_b)$, with serial
number encoding $(d,b)$. This is injective in $b$. No old activity is
overwritten: $K_{s(b)}=0$ before the marked coefficient is installed.

For $v\in F$ let $P_v$ be the primary, masked part of the section, and
let
\[
 f_v(U)=S(\theta^{-1}\mathbin{\cdot_F}v)(o,U)-P_v(U).
 \label{eq:terminal-supplement}
\]
Here $\cdot_F$ is the graded dilation: on the coordinate of a basis
vector $e_j$ of the frame it multiplies by $\theta^{-s_j}$, where $s_j$
is the factor grade of $e_j$, the second entry of its grade $(r_j,s_j)$
in Section~\ref{sec:ex-marked-portfolio-calibration}. The primary part retains the vacuum
coordinate everywhere and the scaled nonvacuum part on healthy
configurations. Hence $f_v=0$ when $Y$ is healthy; otherwise it equals
the scaled nonvacuum section. The map $v\mapsto f_v$ is complex-linear,
measurable, gauge invariant, and depends only on the field and flags in
$Y$. These assertions follow directly from the corresponding properties
of the sound section and the mask. If $q=\max_js_j$ is the maximum
factor grade and each axis section $S(e_j)(o,\cdot)$ has modulus at most
$C_d\ge1$, then
\[
 |f_v(U)|\le C_d\theta^{-q}\|\operatorname{core}v\|_1.
 \tag{T1}
\]
Indeed expand the nonvacuum coefficient in the finite axis basis and
use $0<\theta\le1$, so $\theta^{-s_j}\le\theta^{-q}$ for every axis.

Define the new forced coefficient at $s(b)$ by
\[
 J_v(s(b),U)=\mathbf1_{\{T_b\cap\operatorname{flags}(U)=I_b\}}
                f_v(U)\prod_{i\in b}K_i(U),
 \tag{T2}
\]
and set it to zero at all other addresses. Its support and inventory
are exactly $T_b,I_b$. When it is nonzero, the indicator ensures the
original additional predicate of the reserved address. Old factors
$K_i$ retain all their own predicates. The dependence, invariance and
measurability requirements thus hold on the original dictionary.
The marked coefficient has lead $v$, zero regular correction, and
forced part $J_v$.

\subsection{The exact density and forced profile}
Consider an old compatible collection $c$ at a configuration where
$f_v(U)\ne0$. Split it uniquely into the labels meeting $Y$ and the
labels disjoint from $Y$:
\[
 b=\{i\in c:X_i\cap Y\ne\varnothing\},\qquad r=c\setminus b.
\]
The set $b$ is nonempty. Indeed $Y$ is then not healthy, and a flagged
site of $Y$ lies in some inventory of $c$, hence in its support.
Replace $b$ by $s(b)$ and leave $r$ unchanged. The result is compatible:
its old remaining supports are disjoint from $Y$ and from all supports
in $b$, hence from $T_b$. Inventories still cover exactly the same
flags. Conversely, from a compatible marked collection containing
$s(b)$, replace that address by $b$. Its other supports are disjoint
from $T_b$, so this inverse restores precisely the original collection.
These two operations are inverse before any summation.

For each paired term, (T2) times its remaining old activity factors is
$f_v(U)$ times the original product. The following absolute estimate
justifies summing these identities. Suppose $k_i\ge0$ majorizes
$|K_i|$ and has the actual weighted site bound
\[
 \sum_{i:x\in X_i} M(X_i)k_i\le w,
 \qquad M(X)=\exp\{E\,\operatorname{cover}_B(X)\},
 \tag{T3}
\]
where $E\ge0$ is the purchased support exponent. Because every $X_i$ in
$b$ meets $Y$, the connected covering inequality gives
\[
 M(T_b)\prod_{i\in b}k_i\le M(Y)\prod_{i\in b}M(X_i)k_i.
\]
Moreover the sum of $M(X_i)k_i$ over labels meeting $Y$ is at most
$|Y|w$: sum (T3) over $x\in Y$. For every finite set of such labels,
the sum of products over its nonempty subsets is
$\prod_i(1+M(X_i)k_i)-1\le e^{|Y|w}-1$.
Increasing the finite sets and using nonnegativity proves
\[
 \sum_b M(T_b)\sup_U|J_v(s(b),U)|
 \le C_d\theta^{-q}\|\operatorname{core}v\|_1
             M(Y)(e^{|Y|w}-1).
 \tag{T4}
\]
Dropping the compatibility restrictions only enlarges this bound.
It supplies a summable majorant for the new coefficients. The lawful
old gas bounds supply the remaining dressing majorants, so the
regrouping above is absolutely convergent. Its sum is exactly $f_v$
times the old gas. Adding the primary contribution proves the pointwise
identity
\[
 \operatorname{density}(v)(U)
   =R(U)S(\theta^{-1}\mathbin{\cdot_F}v)(0,U)
   =a_N^d R(U)\,\mathrm{initial}_{N,p}(v)(U).
 \tag{T5}
\]
Here $R(U)$ denotes the value at $U$ of the unmarked density of the
terminal representation, and $a_N$ is the mesh at level $N$. The last
equality is the definition of the physical normalization
$\mathrm{initial}_{N,p}(v)=a_N^{-d}S(\theta^{-1}\mathbin{\cdot_F}v)(0,\cdot)$.

There is also an identity before multiplication by the unmarked gas.
Group the forced coefficients according to their support and inventory.
The address $s(b)$ contributes exactly the expression in (T2), and
all addresses outside its image contribute zero. The bijection between
$b$ and the finite compatible families meeting $Y$ identifies this
profile with the attachment of the singleton response $f_v$ on $Y$.
This is the forced-profile identity needed by the adjacent transfer
construction. It follows from equality of the individual products and
the summable majorant (T4), not from the density identity (T5).

The footprint, graded dilation and healthy mask commute with signed
lattice motions. The sound section also commutes with the independent
complex-conjugation action. Therefore $f_v$ transforms into
$f_{g\cdot v}$, with the same scalar conjugation when appropriate.
Reindexing the compatible collections preserves their supports and
inventories and proves the same covariance for the attached profile.
No equivariant choice of the serial numbers is needed: the profile
sums over addresses. The density, coefficient index and forced profile
thus have the common motion covariance required at terminal birth.

\subsection{Paying the birth cost before choosing the history}
It remains to pay the factor in (T4), uniformly in the period. The
axis bound $C_d$ comes from the fixed finite local section. Its stencil
has cardinality at most a fixed $R_d$, and, once the cover radius
contains that stencil, $\operatorname{cover}_B(Y)\le1$. Thus the fee is
at most
\[
 C_d\theta^{-q}e^E(e^{R_dw}-1).
 \tag{T6}
\]
On a fixed profile stage, $E,C_d,R_d,q$ are fixed. For the hatted
coupling $x$ the chosen threshold satisfies $\theta^{-1}\le x$ once
$x\ge1$, and the forced allowance is
\[
 w(x)=\exp\{-\exp((\log\log x)^2/5)\},\qquad 0<w(x)\le1.
\]
For each fixed integer $n$, $\sqrt{w(x)}x^n\to0$. To see this, put
$u=\log\log x$. The logarithm of that product is
$-\tfrac12e^{u^2/5}+ne^u$, which tends to $-\infty$ because
$e^{u^2/5-u}\to\infty$. Also
$e^z-1\le ze^z$ for $z\ge0$, by integration of $e^t$ over $[0,z]$.
Consequently (T6) is at most
\[
 K_dx^qw(x),\qquad K_d=C_de^ER_de^{R_d}.
\]
Choose $x\ge\max(1,K_d)$ and
$\sqrt{w(x)}x^{q+1}\le1$. Then this last expression is at most
$\sqrt{w(x)}$. This proves the terminal remainder bound with charge
$\|\operatorname{core}v\|_1$, uniformly in $N,p$ at this profile stage.

At stage $j$ there are only finitely many grades and earlier profiles
through $j$. Take a common maximum of their fee thresholds and stencil
fit thresholds, then take running maxima in $j$. Include this delay
among the purchases of Section~\ref{sec:common-observable-purchases}
before applying the common scale selection. At level $m$, both its
grade $d\le J(m)$ and its actual profile index $J(m)$ lie in the
finite list purchased at $J(m+1)$. The fee therefore fits at every
admitted terminal level, on the same selected scale sequence. With
(T1)--(T5), this supplies the terminal marked data and all of their
bounds for the backward transfer construction.

\subsection{Summary of the terminal construction}
For later reference we collect what has been proved, with its exact
inputs. Fix $d$, $N$ and $p$ as above.

\begin{proposition}[Terminal marked coefficients]
\label{prop:terminal-marked-voice}
Assume the following.
\begin{enumerate}
\item The terminal representation $R$ of the selected history at $(N,p)$
  is lawful: its forced labels satisfy the regularity, measurability and
  locality requirements of the dictionary, and their activities have
  majorants $k_i\geq|K_i|$ obeying the weighted site bound (T3) with
  some $w>0$.
\item For every pair $I\subseteq X$ with $I\neq\varnothing$ the
  dictionary of $R$ contains infinitely many reserved addresses with
  support $X$, inventory $I$ and zero activity, whose additional
  predicate holds at every configuration with
  $X\cap\operatorname{flags}(U)=I$.
\item The local section $S$ is sound, and
  $|S(e_j)(0,U)|\leq C_d$ for all axes $j$ and configurations $U$, with
  $C_d\geq1$.
\item The fee inequality
  $C_d\theta^{-q}M(Y)(e^{|Y|w}-1)\leq\sqrt w$ holds.
\end{enumerate}
For $v\in F$ let $\mathcal V(v)$ be the marked coefficient with lead
$v$, zero regular correction and forced part $J_v$ defined in (T2). Then
$v\mapsto\mathcal V(v)$ is complex-linear in each of its three entries,
and for every $v$:
\begin{enumerate}
\item each forced entry $J_v(i,\cdot)$ is gauge invariant and measurable,
  depends only on the field and flags on the support of the address $i$,
  and is nonzero only at configurations where that address is eligible
  for the original dictionary of $R$;
\item the forced entries have majorants with
  $\sum_bM(T_b)\sup_U|J_v(s(b),U)|
     \leq\sqrt w\,\|\operatorname{core}v\|_1$;
\item the density identity (T5) holds at every configuration;
\item the forced profile of $J_v$ on the activities of $R$ equals the
  attachment of the singleton response $f_v$ on the footprint $Y$;
\item for every signed lattice motion, with or without complex
  conjugation, the density of $\mathcal V(v)$ transported by the motion
  is the density of $\mathcal V$ at the transported label.
\end{enumerate}
On the selected scale sequence, hypothesis~(4) holds at every admitted
terminal level $N$, for every grade $d\leq J(N)$ and every period, once
the entrance delays include the fee thresholds chosen above.
\end{proposition}
\begin{proof}
Linearity holds because $v\mapsto f_v$ is linear and (T2) is linear in
$f_v$. Assertion~(1) was verified after (T2); the eligibility for the
original dictionary uses the indicator in (T2) and hypothesis~(2).
Assertion~(2) is (T4) combined with hypothesis~(4). Assertions~(3), (4)
and~(5) are respectively the identity (T5), the forced-profile identity
and the covariance proved in the second subsection. The final sentence is
the conclusion of the third subsection.
\end{proof}

Hypothesis~(2) is a property of the selected history, not of the
construction in this section. The reserved-address property is an input
here; the first subsection only recalls why the dictionary can carry such
addresses, and the realization of the property on the selected history
must be supplied by the history-selection construction. Hypothesis~(1)
is likewise a property of the selected history. The inclusion of the fee
thresholds among the entrance delays, used in the final sentence of the
proposition, is a condition on the schedule which is not discharged
here. These hypotheses are recorded as inputs and are consumed,
together with the five conclusions, by the backward constructions of
Sections~\ref{sec:selected-centered-service} and
\ref{sec:actual-marked-portfolio-calibration}.
